\section{Conservation, maximality, and initial bias}\label{sec:mechanics}
\begin{lemma}[Mechanical identities and global quadratic bounds]
\label{lem:mechanics}
For either microscopic flow, almost surely on its defined trajectories,
\begin{align}
 \pi_t^\eps(|v|^2)&=\pi_0^\eps(|v|^2),\label{eq:energy}\\
 \pi_t^\eps(x\cdot v)&=\pi_0^\eps(x\cdot v)+t\pi_0^\eps(|v|^2)
                                  +\eps\cH_\eps(t),\label{eq:virial}\\
 \pi_t^\eps(|x|^2)&=\pi_0^\eps(|x|^2)+2t\pi_0^\eps(x\cdot v)
       +t^2\pi_0^\eps(|v|^2)+2\eps\int_0^t\cH_\eps(s)\dd s.
 \label{eq:spread}
\end{align}
Mass, momentum, and angular momentum are constant, and the position sum is
its initial value plus $t$ times momentum. For every fixed $q\in\cQ$,
\begin{equation}\label{eq:global-bound}
 \sup_{t\le T}\sum_j|q(z_j(t))|\le C_{q,T}N.
\end{equation}
\end{lemma}
\begin{proof}
Energy conservation and the initial velocity support give
$\|V(t)\|_2\le R_v\sqrt N$. With continuous positions and their piecewise
velocity derivative, Minkowski gives
\[
 \|X(t)\|_2\le\|X(0)\|_2+\int_0^t\|V(s)\|_2\dd s
                           \le(R_x+TR_v)\sqrt N.
\]
Cauchy--Schwarz controls the absolute mixed and angular sums and the linear
position and velocity sums, proving \eqref{eq:global-bound}.
At a contact the pair energy and momentum jumps are zero, whereas
\[
 x\cdot(v'-v)+(x+\eps\omega)\cdot(w'-w)=\eps B.
\]
The angular jump is $\eps B\,\omega\times\omega=0$.
Streaming gives derivatives $|v|^2$ for $x\cdot v$, $2x\cdot v$ for
$|x|^2$, and zero for $x\times v$. Integration proves the identities.
Accepted pasted contacts have exactly the same separation and orientation;
rejected contacts and restarts have no phase impulse. Overlaps introduce
no additional term. The nonnegative cumulative speed sum is bounded by $C_TN/\eps$,
so no extra true collision-count moment estimate is required for its
integrability. H1 justifies taking expectations.
\end{proof}

\begin{proposition}[Maximal transport-closed quadratic invariants]
\label{prop:maximal}
The definition of $\cQ$ gives exactly \eqref{eq:Q}. In spatial dimension
$d$, its dimension is $2d+4+d(d-1)/2$.
\end{proposition}
\begin{proof}
Write a general quadratic as
\[
 \alpha+b\cdot x+c\cdot v+x^TDx+x^TMv+v^TFv,
 \qquad D^T=D,\quad F^T=F.
\]
At a common collision position the linear-in-velocity terms depend only
on pair momentum. Expressing the velocities by their center velocity and
relative velocity $g$, pair invariance of the remaining term requires
$g^TFg$ invariant under every reflection of $g$. A reflection can send a
vector to any other vector of the same length; hence $F=\kappa I$.
The velocity-quadratic coefficient of $q(x+tv,v)$ is
$F+t\operatorname{sym}M+t^2D$. For it to be scalar at every $t$,
$\operatorname{sym}M=dI$ and $D=eI$. The skew part of $M$ is arbitrary.
Conversely these conditions make all the transported tests collision
invariant. They give \eqref{eq:Q}. Counting its independent coefficients
gives the dimension. This proves maximality with the stated property,
without asserting invariance for an arbitrary anisotropic quadratic.
\end{proof}

For $q\in\cQ$, the mechanical identities combine to
\begin{equation}\label{eq:Q-identity}
 \pi_t^\eps(q)=\pi_0^\eps(r_q(t))
 +\eps\left[d(q)\cH_\eps(t)+2e(q)\int_0^t\cH_\eps(s)\dd s\right].
\end{equation}
The skew term in \eqref{eq:Q} has neither a streaming derivative nor a
finite-diameter impulse. Boltzmann's co-located invariants give
\begin{equation}\label{eq:Bol-Q}
 f_t(q)=f_0(r_q(t)).
\end{equation}
These deterministic moment identities can be justified first with spatial
and velocity cutoffs. Global energy and the propagated spatial second
moment remove them; the latter follows from compact initial support and
$2|x\cdot v|\le|x|^2+|v|^2$. This deterministic integration is not an
unbounded fluctuation theorem.

\begin{lemma}[Initial vacancy bias]\label{lem:vacancy}
For every fixed polynomial $q$,
\begin{equation}\label{eq:vacancy}
 a_\eps(q):=\eps^{-1}[\E\pi_0^\eps(q)-f_0(q)]\longrightarrow a(q).
\end{equation}
\end{lemma}
\begin{proof}
The insertion identity for \eqref{eq:ensemble} is
$F_{1,0}^\eps(x,v)=f_0(x,v)\Prob\{K_x=0\}$, where $K_x$ counts centers
of an independent initial configuration in the radius-$\eps$ ball about $x$.
H1, smoothness of $\rho_0$, and $\mu\eps^3=\eps$ give, uniformly in $x$,
\[
 \E K_x=\eps\bex\rho_0(x)+O(\eps^2),\qquad
 \E K_x(K_x-1)\le\mu^2\left(\int_{B(x,\eps)}\rho_0\right)^2
                                                   =O(\eps^2).
\]
The elementary Bonferroni inequalities
\[
 \E K_x-\tfrac12\E K_x(K_x-1)
                 \le\Prob(K_x\ge1)\le\E K_x
\]
yield
$F_{1,0}^\eps=f_0[1-\eps\bex\rho_0+O(\eps^2)]$.
Multiply by $q$ and integrate on the compact support of $f_0$.
This proves \eqref{eq:vacancy}. In particular the activity $\mu$ cannot
be identified with $\E N$ at the order relevant to Boltzmann centering.
\end{proof}

\begin{lemma}[True mean transfer]\label{lem:mean-transfer}
For $q\in\cQ$ and every fixed integer $p>21$,
\begin{equation}\label{eq:mean-transfer}
 \sup_{t\le T}\eps^{-1}
  |\E\pi_t^\eps(q)-\E\pi_t^{\eps,p}(q)|
                      \le C_{q,T,p}\eps^{(p-21)/(20p)}.
\end{equation}
Moreover $\E\cH_\eps\to H$ uniformly.
\end{lemma}
\begin{proof}
Let $E_\eps$ be H6's full-time discrepancy event. Off it the moment paths
agree; on it \eqref{eq:global-bound} bounds their absolute particle-sum
difference by $C_{q,T}N$. Holder and H1 give
\[
 \eps^{-1}\E[(N/\mu)\ind_{E_\eps}]
 \le C_p\eps^{-1+(21/20)(1-1/p)}
 =C_p\eps^{(p-21)/(20p)}.
\]
For $p=42$ the exponent is $1/40$. Apply \eqref{eq:virial} to both
coupled flows: their initial and energy terms agree, so
$\cH_\eps-\cH_\eps^p=\eps^{-1}
[\pi_t^\eps(x\cdot v)-\pi_t^{\eps,p}(x\cdot v)]$.
The mean comparison therefore controls the speed-sum expectation uniformly.
Lemma~\ref{lem:incidence-mean} completes the proof. This is first-moment
transfer, not a deduction from probability coupling alone.
\end{proof}
